% GENERATED FILE. Edit canonical lessons, then run npm run generate:books.
% canonical-source-hash: fnv1a64:225f07ec0aec8d85
% canonical-lessons: SA-C105-duradarsanam, SA-C105-sanganakam, SA-C105-yantram, SA-C105-vyajanam, SA-C105-petika

\chapter{Television, Computer, Machine, Fan, Box}
\label{ch:sa-television-computer-machine-fan-box}
\hlchaptermodality{105}

\begin{chapteropening}
\textbf{By the end of this chapter:} \emph{I can say a television, a computer, a machine, a fan and a box.}

Complete the last lesson of chapter 105: i can say a television, a computer, a machine, a fan and a box.
\end{chapteropening}

\section[\texorpdfstring{\sk{दूरदर्शनम्} (dūradarśanam)}{दूरदर्शनम् (dūradarśanam)}]{\sk{दूरदर्शनम्} (dūradarśanam) — a television}

\label{lesson:SA-C105-duradarsanam}



\begin{quote}
Before the new one: say the Sanskrit for a shirt, then the Sanskrit for a telephone.
\end{quote}

\subsection*{You'll want to know: \sk{दूरदर्शनम्}}
\textbf{\sk{दूरदर्शनम्}} — \emph{dūradarśanam} — \textquotedblleft{}a television\textquotedblright{}.

\begin{grammarlens}[title={What you've built}]
One more word for this chapter.
\end{grammarlens}

\subsection*{Your turn}
\begin{itemize}
\raggedright
  \item \emph{Say it:} \emph{dūradarśanam}
  \item \emph{Say it:} \emph{dūradarśanam}, once more
  \item \emph{Say it:} say \emph{dūradarśanam}
  \item \emph{Recall:} say the Sanskrit for a shirt, then the Sanskrit for a telephone, then say \emph{dūradarśanam} again
\end{itemize}

\begin{tcolorbox}[breakable,colback=teal!4,colframe=teal!35!black,title={Before you move on}]
What does \sk{दूरदर्शनम्} mean? (\textquotedblleft{}A television\textquotedblright{}.) Say it once more.
\end{tcolorbox}
\section[\texorpdfstring{\sk{संगणकम्} (saṅgaṇakam)}{संगणकम् (saṅgaṇakam)}]{\sk{संगणकम्} (saṅgaṇakam) — a computer}

\label{lesson:SA-C105-sanganakam}



\begin{quote}
Before the new one: say the Sanskrit for a telephone, then the Sanskrit for a television.
\end{quote}

\subsection*{You'll want to know: \sk{संगणकम्}}
\textbf{\sk{संगणकम्}} — \emph{saṅgaṇakam} — \textquotedblleft{}a computer\textquotedblright{}.

\begin{grammarlens}[title={What you've built}]
One more word for this chapter.
\end{grammarlens}

\subsection*{Your turn}
\begin{itemize}
\raggedright
  \item \emph{Say it:} \emph{saṅgaṇakam}
  \item \emph{Say it:} \emph{saṅgaṇakam}, once more
  \item \emph{Say it:} say \emph{saṅgaṇakam}
  \item \emph{Recall:} say the Sanskrit for a telephone, then the Sanskrit for a television, then say \emph{saṅgaṇakam} again
\end{itemize}

\begin{tcolorbox}[breakable,colback=teal!4,colframe=teal!35!black,title={Before you move on}]
What does \sk{संगणकम्} mean? (\textquotedblleft{}A computer\textquotedblright{}.) Say it once more.
\end{tcolorbox}
\section[\texorpdfstring{\sk{यन्त्रम्} (yantram)}{यन्त्रम् (yantram)}]{\sk{यन्त्रम्} (yantram) — a machine}

\label{lesson:SA-C105-yantram}



\begin{quote}
Before the new one: say the Sanskrit for a television, then the Sanskrit for a computer.
\end{quote}

\subsection*{You'll want to know: \sk{यन्त्रम्}}
\textbf{\sk{यन्त्रम्}} — \emph{yantram} — \textquotedblleft{}a machine\textquotedblright{}.

\begin{grammarlens}[title={What you've built}]
One more word for this chapter.
\end{grammarlens}

\subsection*{Your turn}
\begin{itemize}
\raggedright
  \item \emph{Say it:} \emph{yantram}
  \item \emph{Say it:} \emph{yantram}, once more
  \item \emph{Say it:} say \emph{yantram}
  \item \emph{Recall:} say the Sanskrit for a television, then the Sanskrit for a computer, then say \emph{yantram} again
\end{itemize}

\begin{tcolorbox}[breakable,colback=teal!4,colframe=teal!35!black,title={Before you move on}]
What does \sk{यन्त्रम्} mean? (\textquotedblleft{}A machine\textquotedblright{}.) Say it once more.
\end{tcolorbox}
\section[\texorpdfstring{\sk{व्यजनम्} (vyajanam)}{व्यजनम् (vyajanam)}]{\sk{व्यजनम्} (vyajanam) — a fan}

\label{lesson:SA-C105-vyajanam}



\begin{quote}
Before the new one: say the Sanskrit for a computer, then the Sanskrit for a machine.
\end{quote}

\subsection*{You'll want to know: \sk{व्यजनम्}}
\textbf{\sk{व्यजनम्}} — \emph{vyajanam} — \textquotedblleft{}a fan\textquotedblright{}.

\begin{grammarlens}[title={What you've built}]
One more word for this chapter.
\end{grammarlens}

\subsection*{Your turn}
\begin{itemize}
\raggedright
  \item \emph{Say it:} \emph{vyajanam}
  \item \emph{Say it:} \emph{vyajanam}, once more
  \item \emph{Say it:} say \emph{vyajanam}
  \item \emph{Recall:} say the Sanskrit for a computer, then the Sanskrit for a machine, then say \emph{vyajanam} again
\end{itemize}

\begin{tcolorbox}[breakable,colback=teal!4,colframe=teal!35!black,title={Before you move on}]
What does \sk{व्यजनम्} mean? (\textquotedblleft{}A fan\textquotedblright{}.) Say it once more.
\end{tcolorbox}
\section[\texorpdfstring{\sk{पेटिका} (peṭikā)}{पेटिका (peṭikā)}]{\sk{पेटिका} (peṭikā) — a box}

\label{lesson:SA-C105-petika}



\begin{quote}
Before the new one: say the Sanskrit for a machine, then the Sanskrit for a fan.
\end{quote}

\subsection*{You'll want to know: \sk{पेटिका}}
\textbf{\sk{पेटिका}} — \emph{peṭikā} — \textquotedblleft{}a box\textquotedblright{}.

\begin{grammarlens}[title={What you've built}]
One more word for this chapter.
\end{grammarlens}

\subsection*{Your turn}
\begin{itemize}
\raggedright
  \item \emph{Say it:} \emph{peṭikā}
  \item \emph{Say it:} \emph{peṭikā}, once more
  \item \emph{Say it:} say \emph{peṭikā}
  \item \emph{Recall:} say the Sanskrit for a machine, then the Sanskrit for a fan, then say \emph{peṭikā} again
\end{itemize}

\begin{tcolorbox}[breakable,colback=teal!4,colframe=teal!35!black,title={Before you move on}]
What does \sk{पेटिका} mean? (\textquotedblleft{}A box\textquotedblright{}.) Say it once more.
\end{tcolorbox}
